\section{Datensouveränität und Systemarchitektur} \label{sec:inhouse}

\subsection{Rechtliche Anforderungen im Gesundheitswesen} \label{sec:anforderungen}

Medizinische Bilddaten zählen nach Artikel~9 der Datenschutz-Grundverordnung (DSGVO) zu
den besonderen Kategorien personenbezogener Daten. Ihre Verarbeitung erfordert eine
explizite Rechtsgrundlage sowie organisatorische und technische Schutzmaßnahmen
\cite{gdprMedical}. Ergänzend klassifiziert der EU AI Act -- seit August 2024 in Kraft --
KI-gestützte Diagnosetools in Geräten der Klasse~IIb (z.\,B.\ Röntgen- oder CT-Systeme)
als Hochrisiko-KI-Systeme \cite{euAiAct2024}. Hersteller und Betreiber müssen bis
August 2026 vollständige Konformität nachweisen, einschließlich technischer Dokumentation,
Transparenzanforderungen und Risikoklassifikation. Die Wahl der Systemarchitektur ist
damit zugleich eine regulatorische Entscheidung.

\subsection{Cloud vs.\ On-Premise} \label{sec:cloud_onpremise}

Kommerzielle Cloud-Infrastruktur bietet Vorteile hinsichtlich Rechenkapazität und
Skalierbarkeit. Jedoch ist eine geografisch \glqq EU-Region\grqq{} bei Anbietern wie
AWS, Azure oder Google Cloud keine ausreichende Garantie für Datensouveränität: Der
US CLOUD Act verpflichtet US-amerikanische Unternehmen, behördlichen Zugriffsanfragen
auch aus europäischen Rechenzentren nachzukommen \cite{gdprMedical}. Ein faktischer
Datentransfer in Drittstaaten ist damit möglich -- und mit den Anforderungen des
Artikel~9 DSGVO nur schwer vereinbar.

On-Premise-Architekturen schließen dieses Risiko strukturell aus: Patientendaten
verbleiben vollständig innerhalb der Einrichtung, die Nachweispflichten der DSGVO sind
leichter erfüllbar, und die Hochrisikoanforderungen des EU AI Act lassen sich
kontrollierter umsetzen. In der industriellen Praxis wurde dieser Zielkonflikt durch
Inhouse-Inferenzsysteme mit Edge-AI-Ansätzen adressiert -- eine Architektur, die eine
analoge Umsetzung im Krankenhauskontext ermöglicht.

\subsection{Federated Learning als datenschutzkonformer Lernverbund} \label{sec:federated}

KI-Modelle profitieren von größeren Trainingsdatensätzen -- eine zentrale Aggregation
von Patientendaten über Einrichtungsgrenzen hinweg ist jedoch datenschutzrechtlich nicht
zulässig. Federated Learning adressiert diesen Zielkonflikt: Jede Einrichtung trainiert
ein lokales Modell auf eigenen Daten; lediglich die aggregierten Modellgewichte werden
geteilt. Rohdaten verbleiben vollständig innerhalb der jeweiligen Institution
\cite{federatedLearning2022}. Ergänzende Techniken wie Differential Privacy -- durch
statistisches Rauschen in den Gradienten -- und Secure Multi-Party Computation stärken
den Datenschutz weiter. Aktuelle Forschungsergebnisse belegen, dass
Federated-Learning-Modelle bei ausreichend großem Klinikverbund eine mit zentral
trainierten Modellen vergleichbare Qualität erreichen \cite{federatedAdvances2025}.
